% Part: turing-machines
% Chapter: undecidability
% Section: representing-tms

\documentclass[../../../include/open-logic-section]{subfiles}

\begin{document}

\olfileid{tur}{und}{rep}
\olsection{Nggambarake Mesin Turing}

\begin{explain}
Kanggo nggambarake mesin Turing lan tumindake nganggo
!!a{sentence} saka logika orde kapisan, kita kudu
netepake basa kang cocog. Basa iki dumadi saka rong
bagean: !!{predicate}s kanggo njlentrehake konfigurasi
mesin, lan ekspresi kanggo menehi nomer marang
langkah-langkah lumakune mesin (``wayah'') lan
panggonan ing pita.

Kita ngenalake rong jinis !!{predicate}s kang
padha-padha nduwé rong papan: kanggo saben
kahanan~$q$, !!a{predicate}~$\Obj Q_q$,
lan kanggo saben simbol pita~$\sigma$,
!!a{predicate}~$\Obj S_\sigma$. Kang kapisan
ngidini kita njlentrehake kahanan~$M$ lan
panggonan sirah maca-nulis, dene kang kapindho
ngidini kita njlentrehake isi pita.

Kanggo nyatakake panggonan sirah maca-nulis
lan cacahe langkah kang wis ditindakake,
kita mbutuhake cara kanggo nyatakake wilangan.
Kanggo iki kita nganggo !!a{constant}~$\Obj 0$
lan fungsi $1$-papan~$\prime$, yaiku fungsi
penerus. Miturut pranatan, fungsi iki ditulis
\emph{sawise} argumene (tanpa tandha kurung).

Kanggo saben wilangan $n$ ana suku baku~$\num{n}$,
yaiku \emph{numeral} kanggo~$n$, kang
nggambarake wilangan iku ing~$\Lang L_M$.
$\num{0}$ yaiku $\Obj 0$, $\num{1}$
yaiku $\Obj 0'$, $\num{2}$ yaiku $\Obj 0''$,
lan sateruse. Sacara luwih formal:
\begin{align*}
\num{0} & = \Obj 0 \\
\num{n+1} &= \num{n}'
\end{align*}
Suku $\num{0}$, yaiku $\Obj 0$, nyebut
panggonan paling kiwa ing pita lan wektu
sadurunge langkah lumaku kapisan (konfigurasi
wiwitan). Suku $\num{1}$, yaiku $\Obj 0'$,
nyebut kothak ing sisih tengen kothak paling
kiwa lan wektu sawisé langkah kapisan,
lan sateruse.

Kita uga ngenalake !!a{predicate}~$<$ kanggo
nyatakake urutan panggonan ing pita (kanthi
teges ``ing sisih kiwa'') lan urutan langkah
lumaku (kanthi teges ``sadurunge'').

Sawisé basane ditetepake, kita dhaptar
``aksioma'' saka $!T(M, w)$, yaiku
!!{sentence}s kang bebarengan njlentrehake
tumindake~$M$ nalika diwenehi input~$w$.
Ana !!{sentence}s kang netepake syarat
kanggo $\Obj 0$, $\prime$, lan $<$,
!!{sentence}s kang njlentrehake konfigurasi
input, lan !!{sentence}s kang njlentrehake
konfigurasi $M$ sawisé nindakake instruksi
tartamtu.
\end{explain}

\begin{defn}
  \ollabel{defn:tm-descr}
Kanggo mesin Turing $M = \tuple{Q, \Sigma, q_0, \delta}$,
basa~$\Lang L_M$ dumadi saka:
\begin{enumerate}
\item !!{predicate} rong papan $\Obj Q_q(x, y)$
  kanggo saben kahanan~$q \in Q$. Miturut intuisi,
  $\Obj Q_q(\num{m}, \num{n})$ nyatakake
  ``sawisé $n$ langkah, $M$ ana ing kahanan~$q$
  lan lagi maca kothak kapangka-$m$.''
\item !!{predicate} rong papan $\Obj S_\sigma(x, y)$
  kanggo saben simbol~$\sigma\in \Sigma$.
  Miturut intuisi, $\Obj S_\sigma(\num{m},
  \num{n})$ nyatakake ``sawisé $n$ langkah,
  kothak kapangka-$m$ ngemot simbol~$\sigma$.''
\item !!{constant} $\Obj 0$
\item !!{function} siji papan $\prime$
\item !!{predicate} rong papan $<$
\end{enumerate}
\end{defn}

!!{sentence}s kang njlentrehake lumakune mesin Turing~$M$
kanthi input $w = \sigma_{i_1}\dots\sigma_{i_k}$
yaiku ing ngisor iki:
\begin{enumerate}
\item Aksioma kang njlentrehake wilangan lan~$<$:
\begin{enumerate}
\item !!^a{sentence} kang nyatakake yen saben wilangan
luwih cilik tinimbang peneruse:
\[
\lforall[x][x < x']
\]
\item !!^a{sentence} kang mesthekake yen $<$ transitif:
\[
\lforall[x][\lforall[y][\lforall[z][
      ((x < y \land y < z) \lif x < z)]]]
\]
\end{enumerate}
\item Aksioma kang njlentrehake konfigurasi input:
\begin{enumerate}
\item Sawisé $0$~langkah---sadurunge mesin miwiti
  pakaryane---$M$ ana ing kahanan wiwitan~$q_0$
  lan lagi maca kothak~$1$:
\[
\Obj Q_{q_0}(\num{1}, \num{0})
\]
% OLPL-155: Kanggo k=0, runtutan simbol input kosong lan suku sigma_i1 ing tampilan iki butuh pranatan konjungsi kosong.
\item Kothak-kothak $k+1$ kapisan ngemot simbol
  $\TMendtape$, $\sigma_{i_1}$, \dots, $\sigma_{i_k}$:
\[
\Obj S_\TMendtape(\num{0}, \num{0}) \land
\Obj S_{\sigma_{i_1}}(\num{1}, \num{0}) \land
\dots \land
\Obj S_{\sigma_{i_k}}(\num{k}, \num{0})
\]
\item Ing kothak liyane, pita kosong:
\[
\lforall[x][(\num{k} < x \lif \Obj S_\TMblank(x, \num{0}))]
\]
\end{enumerate}
\item Aksioma kang njlentrehake transisi saka
  sawijining konfigurasi menyang konfigurasi sabanjure:

% OLPL-154: Sumber nyebut klausa iki ukara, nanging x lan y tetep variabel bebas ing A(x,y); iki rumus mbukak.
Kanggo bagean sabanjure, $!A(x, y)$ yaiku
konjungsi kabeh !!{sentence}s kanthi wujud
\[
\lforall[z][
  (((z < x \lor x < z) \land \Obj S_\sigma(z, y))
  \lif \Obj S_\sigma(z, y'))]
\]
ing ngendi $\sigma \in \Sigma$.
Kita nganggo $!A(\num{m},\num{n})$ kanggo
nyatakake ``kajaba ing kothak~$m$, isi pita
sawisé $n+1$ langkah padha karo isi pita
sawisé $n$ langkah.''
\begin{enumerate}
\item \ollabel{rep-right} Kanggo saben instruksi $\delta(q_i, \sigma) =
  \tuple{q_j, \sigma', \TMright}$, !!{sentence} iki:
\begin{align*}
& \lforall[x][\lforall[y][(
   (\Obj Q_{q_i}(x, y) \land \Obj S_{\sigma}(x, y)) \lif {}]] \\
&\qquad   (\Obj Q_{q_j}(x', y') \land \Obj S_{\sigma'}(x, y') \land
!A(x, y)))
\end{align*}
Iki nyatakake yen sawisé~$y$ langkah, mesin ana
ing kahanan~$q_i$ lan lagi maca kothak~$x$
kang ngemot simbol~$\sigma$, mula sawisé
$y+1$ langkah mesin maca kothak~$x+1$,
ana ing kahanan~$q_j$, kothak~$x$ saiki
ngemot~$\sigma'$, lan saben kothak saliyane~$x$
ngemot simbol kang padha kaya sawisé~$y$ langkah.

\item \ollabel{rep-left} Kanggo saben instruksi $\delta(q_i, \sigma) =
  \tuple{q_j, \sigma', \TMleft}$, !!{sentence} iki:
\begin{align*}
& \lforall[x][\lforall[y][
    ((\Obj Q_{q_i}(x', y) \land \Obj S_{\sigma}(x', y)) \lif {}]]\\
& \qquad   (\Obj Q_{q_j}(x, y') \land \Obj S_{\sigma'}(x', y') \land
!A(x, y))) \land {}\\
& \lforall[y][((\Obj Q_{q_i}(\num{0}, y) \land \Obj S_{\sigma}(\num{0},
    y)) \lif {}]\\
& \qquad (\Obj Q_{q_j}(\num{0}, y') \land \Obj S_{\sigma'}(\num{0},
  y') \land !A(\num{0}, y)))
\end{align*}
Gatekna sedhela carane iki lumaku: saiki kita
ora miwiti nganggo ``yen maca kothak~$x$ \dots'',
nanging nganggo ``yen maca kothak $x+1$ \dots''.
Obah menyang kiwa ateges ing langkah sabanjure
mesin maca kothak~$x$. Nanging kothak kang
ditulisi yaiku~$x+1$. Kita nindakake iki
amarga ora nduwé operasi pangurangan utawa
fungsi kang menehi wilangan sadurunge.

Elinga yen wilangan kang wujude $x+1$ yaiku
$1$, $2$, \dots; mula iki durung nyakup
kahanan nalika mesin maca kothak~$0$ lan
kudu obah menyang kiwa (kang mesthi ora
bisa---mesin mung tetep ing panggonane).
Kasus mligi iki dicakup dening konjungsi
kapindho: yen sawisé $y$ langkah mesin
maca kothak~$0$ ing kahanan $q_i$ lan
kothak~$0$ ngemot simbol~$\sigma$,
mula sawisé $y+1$ langkah mesin isih
maca kothak~$0$, saiki ana ing kahanan~$q_j$,
simbol ing kothak~$0$ yaiku $\sigma'$,
lan kothak-kothak saliyane kothak~$0$
ngemot simbol kang padha kaya sawisé
$y$~langkah.
\item \ollabel{rep-stay} Kanggo saben instruksi $\delta(q_i, \sigma) =
  \tuple{q_j, \sigma', \TMstay}$, !!{sentence} iki:
\begin{align*}
& \lforall[x][\lforall[y][(
   (\Obj Q_{q_i}(x, y) \land \Obj S_{\sigma}(x, y)) \lif {}]] \\
&\qquad   (\Obj Q_{q_j}(x, y') \land \Obj S_{\sigma'}(x, y') \land
!A(x, y)))
\end{align*}
\end{enumerate}
\end{enumerate}
Tetepna $!T(M, w)$ minangka konjungsi kabeh
!!{sentence}s ing ndhuwur kanggo mesin Turing~$M$
lan input~$w$.

Kanggo nyatakake yen~$M$ pungkasane mandheg,
kita kudu nemokake !!{sentence} kang nyatakake
``sawisé sawetara langkah, fungsi transisi
ora ditegesi.'' Tetepna $X$~minangka himpunan
kabeh pasangan $\tuple{q, \sigma}$ kang
$\delta(q, \sigma)$ ora ditegesi.
Banjur $!E(M, w)$ yaiku !!{sentence}
\[
\lexists[x][\lexists[y][(\bigvee_{\tuple{q, \sigma} \in
      X}(\Obj Q_q(x, y) \land \Obj S_\sigma(x, y)))]]
\]

Yen kita nganggo mesin Turing kang nduwé
kahanan mandheg mligi~$h$, iki malah luwih
gampang: !!{sentence}~$!E(M, w)$
\[
\lexists[x][\lexists[y][\Obj Q_h(x, y)]]
\]
nyatakake yen mesin pungkasane mandheg.

\begin{prop}
\ollabel{prop:mlessk}
Yen $m < k$, mula $!T(M, w) \Entails \num{m} < \num{k}$
\end{prop}

\begin{proof}
Latihan.
\end{proof}

\begin{prob}
Buktekna \olref[tur][und][rep]{prop:mlessk}.
(Pituduh: gunakna induksi marang $k-m$).
\end{prob}

\end{document}
